\section*{अध्याय \ref{ch:g6:sun-earth-moon} --- सूर्य--पृथ्वी--चंद्रमा तंत्र}

\begin{solution}{exo:g6:sun-earth-moon:1}
पृथ्वी अपने \omterm{def:g6:sun-earth-moon:axisorbit}{अक्ष} पर घूमती है --- दिन। पृथ्वी सूर्य के चारों ओर \omterm{def:g6:sun-earth-moon:axisorbit}{कक्षा}
में चलती है --- \omterm{def:g3:sun-days-seasons:year}{वर्ष}। \omterm{def:g2:moon-in-the-sky:moon}{चंद्रमा} पृथ्वी के चारों ओर \omterm{def:g6:sun-earth-moon:axisorbit}{कक्षा} में चलता है ---
महीना।
\end{solution}

\begin{solution}{exo:g6:sun-earth-moon:2}
ध्रुवों से गुज़रती वह काल्पनिक रेखा जिसके चारों ओर पृथ्वी घूमती है। वह
लगभग समकोण के चौथाई जितनी झुकी है --- और पूरी \omterm{def:g6:sun-earth-moon:axisorbit}{कक्षा} में एक ही ओर झुकी
रहती है: झुकाव और वफ़ादारी, दोनों मिलकर ऋतुएँ बनाते हैं।
\end{solution}

\begin{solution}{exo:g6:sun-earth-moon:3}
पृथ्वी के एक चक्कर में $365$ दिन और एक चौथाई लगता है; चार बचाई गई
चौथाइयाँ मिलकर अधिवर्ष का अतिरिक्त दिन बनाती हैं। उसके बिना पंचांग
\omterm{def:g3:sun-days-seasons:year}{ऋतुओं} में से खिसकता चला जाता --- हर चार साल में लगभग एक दिन, और सौ साल
में कई हफ़्ते।
\end{solution}

\begin{solution}{exo:g6:sun-earth-moon:4}
\omterm{def:g2:moon-in-the-sky:moon}{चंद्रमा} तक लगभग $30$ पृथ्वियाँ; और सूर्य के चेहरे पर लगभग $109$
पृथ्वियाँ।
\end{solution}

\begin{solution}{exo:g6:sun-earth-moon:5}
सींक \omterm{def:g6:sun-earth-moon:axisorbit}{अक्ष} की भूमिका निभाती है, और असली \omterm{def:g6:sun-earth-moon:axisorbit}{अक्ष} तारों के मुक़ाबले अपना
झुकाव पक्का बनाए रखता है। उसे घूमने दो --- मान लो वह हमेशा लैंप की ओर
ही झुकी रहे --- तो नमूना कभी न ख़त्म होने वाला ग्रीष्म दिखाएगा:
वफ़ादारी के साथ ऋतुएँ भी ग़ायब हो जाती हैं।
\end{solution}

\begin{solution}{exo:g6:sun-earth-moon:6}
उलटी वाली: जब हमारा गोलार्ध सूर्य की ओर झुकता है, तब दूसरे को दूर
झुकना ही पड़ता है --- एक ही झुकाव दोनों सिरों का पक्ष एक साथ नहीं ले
सकता।
\end{solution}

\begin{solution}{exo:g6:sun-earth-moon:7}
सीधी टॉर्च का \omterm{def:g1:light-and-shadows:source}{प्रकाश} खड़ा और सिमटा हुआ पड़ता है --- छोटा, चमकीला,
काटता हुआ धब्बा। तिरछी करने पर वही \omterm{def:g1:light-and-shadows:source}{प्रकाश} दीवार के लंबे हिस्से पर
पतला और कमज़ोर फैल जाता है। सूर्य की ओर झुकना हमें खड़े और सिमटे हुए
\omterm{def:g1:light-and-shadows:source}{प्रकाश} वाली क़िस्म देता है।
\end{solution}

\begin{solution}{exo:g6:sun-earth-moon:8}
\omterm{def:g2:moon-in-the-sky:moon}{चंद्रमा} की पृथ्वी के चारों ओर अपनी \omterm{def:g6:sun-earth-moon:axisorbit}{कक्षा}: चक्कर लगाते हुए हम उसके
इकलौते धूप वाले आधे हिस्से को बदलते कोणों से देखते हैं --- \omterm{ex:g2:moon-in-the-sky:shapes}{बालचंद्र} से
\omterm{ex:g2:moon-in-the-sky:shapes}{पूर्ण चंद्र} और फिर वापस। कलाओं का तमाशा तीसरी गति चलाती है।
\end{solution}

\begin{solution}{exo:g6:sun-earth-moon:9}
पहला: उत्तरी शिशिर तब पड़ता है जब पृथ्वी सूर्य के ज़रा \emph{ज़्यादा
पास} होती है --- यानी दूरी उलटी दिशा में इशारा करती है। दूसरा: एक ही
पल और एक ही दूरी पर दोनों गोलार्धों में उलटी ऋतुएँ चलती हैं --- कोई
साझा दूरी ऐसा नहीं कर सकती; सिर्फ़ झुकाव कर सकता है, जो एक सिरे का
पक्ष लेता है और दूसरे का नहीं।
\end{solution}

\begin{solution}{exo:g6:sun-earth-moon:10}
झुकाव न हो तो चक्कर का हर पड़ाव एक जैसा दिखेगा: सूर्य साल भर उसी
दोपहर-ऊँचाई तक चढ़ेगा, दिन और रात हर जगह बराबर रहेंगे, और \omterm{def:g3:sun-days-seasons:year}{वर्ष} अपनी
ऋतुएँ खो देगा --- एक लंबा, अनंत वसंत, \omterm{def:g3:sun-days-seasons:sundial}{धूपघड़ी} की चाप जमी हुई, और
पंचांग का सबसे बड़ा नाटक रद्द।
\end{solution}

\begin{solution}{exo:g6:sun-earth-moon:11}
\omterm{def:g2:moon-in-the-sky:moon}{चंद्रमा} जब अपने $27$ दिन का चक्कर लगा रहा था, तब पृथ्वी अपनी \omterm{def:g6:sun-earth-moon:axisorbit}{कक्षा} पर
आगे बढ़ चुकी थी --- इसलिए सूर्य की दिशा खिसक गई, और उसी क़तार तक फिर
पहुँचने के लिए, जो \omterm{ex:g2:moon-in-the-sky:shapes}{पूर्ण चंद्र} बनाती है, \omterm{def:g2:moon-in-the-sky:moon}{चंद्रमा} को लगभग ढाई दिन और
चलना पड़ता है। चक्कर और क़तार अलग-अलग फ़ीते हैं।
\end{solution}

\begin{solution}{exo:g6:sun-earth-moon:12}
ध्रुवीय क़स्बे का रोज़ का वृत्त छोटा होता है और जून के झुकाव की वजह
से वह \emph{पूरा का पूरा} उजाले वाले आधे हिस्से में पड़ सकता है:
क़स्बा घूमता रहता है और कभी दिन से बाहर निकलता ही नहीं --- आधी रात का
सूर्य। दिसंबर में वही झुकाव उसके नन्हे वृत्त को पूरा का पूरा अँधेरे
आधे हिस्से में झुका देता है: पूरे-पूरे चक्कर बिना किसी \omterm{def:g1:day-and-night:daynight}{सूर्योदय} के।
जो झुकाव ध्रुवों पर पूरी तरह करता है, वही हर अक्षांश पर हल्के-हल्के
करता है।
\end{solution}

\begin{solution}{pb:g6:sun-earth-moon:1}
\textbf{1.} हर \omterm{def:g2:measuring-length:units}{सेंटीमीटर} \qty{12800}{km} का।
\textbf{2.} लगभग \qty{0.25}{cm} --- काली मिर्च का दाना (या कोई और
दाना)।
\textbf{3.} गोली से $30 \times 1 = \qty{30}{cm}$ दूर।
\textbf{4.} $109$ \omterm{def:g2:measuring-length:units}{सेंटीमीटर} --- यानी लगभग एक \omterm{def:g2:measuring-length:units}{मीटर}: कोई हूला घेरा या
बड़ी समुद्र-तट वाली गेंद।
\textbf{5.} $150000000 \div 12800 \approx \qty{11700}{cm}$ ---
यानी लगभग \qty{117}{m}।
\textbf{6.} बस ज़रा-सा नहीं: सूर्य वाली गेंद \qty{100}{m} के मैदान की
दूर वाली रेखा से भी लगभग \qty{17}{m} आगे खड़ी करनी पड़ेगी।
\textbf{7.} लगभग $117$ सेकंड --- यानी क़रीब दो मिनट।
\textbf{8.} सूर्य: $109 \div 11700 \approx 0.009$; चंद्रमा-दाना:
$0.25 \div 30 \approx 0.008$ --- लगभग बराबर: गोली से देखने पर दोनों
एक ही नाप के लगते हैं। हमारे आकाश में भी ऐसा ही है --- और जब \omterm{def:g2:moon-in-the-sky:moon}{चंद्रमा}
ठीक सूर्य के सामने से गुज़रता है, तो वह उसे लगभग पूरा ढक लेता है: पूर्ण
सूर्यग्रहण, यानी प्रकृति इसी संयोग का फ़ायदा उठाती हुई।
\textbf{9.} गोली का चक्कर: लगभग $365$ सेकंड --- यानी छह मिनट की सैर;
और दाने का गोली के चारों ओर चक्कर: लगभग $27$ सेकंड।
\textbf{10.} हर सेकंड में एक पूरा घूर्णन --- और बिना किसी निशान वाली
चिकनी एक \omterm{def:g2:measuring-length:units}{सेंटीमीटर} की गोली पर घूर्णन आँख के पकड़ने लायक कुछ छोड़ता ही
नहीं: उस पर एक बिंदु रंग देना ईमानदार हल है।
\textbf{11.} सीधे खड़े से लगभग समकोण के चौथाई जितनी झुकी हुई --- और
पूरे चक्कर में हमेशा मैदान के उसी दूर वाले कोने की ओर ताके हुए, कभी
सूर्य वाली गेंद की ओर मुड़े बिना।
\textbf{12.} जैसे: नमूना आकार और दूरियाँ \emph{साथ-साथ} सच-सच दिखा
देता है --- वह विशाल ख़ालीपन ही उसकी जीत है। जो वह एक साथ नहीं दिखा
सकता वह है सच्ची रफ़्तार वाली गतियाँ (छह मिनट का ``\omterm{def:g3:sun-days-seasons:year}{वर्ष}'' पहले ही
बेतहाशा तेज़ है; असली गोली तो रेंगती) --- और बेशक \omterm{def:g1:light-and-shadows:source}{प्रकाश} तथा गरमी भी:
एक \omterm{def:g2:measuring-length:units}{मीटर} चौड़ी कोई गेंद खेल के मैदान को उस तरह उजाला नहीं देती जिस तरह
सूर्य किसी \omterm{def:g4:solar-system:planet}{ग्रह} को देता है।
\end{solution}
